% AVL rotace: pravá rotace (y,x,A,B,C) <-> levá rotace

\begin{figure}[H]
\centering
\begin{tikzpicture}[
    scale=1.15,
    every node/.style={
        circle,
        draw,
        fill=white,
        inner sep=0pt,
        minimum size=22pt,
        font=\Large\itshape
    },
    line width=0.9pt
]

% ============================================================
% GLOBÁLNÍ PARAMETRY
% ============================================================

\def\SiblingDist{1.15}  % vzdálenost sourozeneckých uzlů/trojúhelníků
\def\LevelDist{1.1}     % vertikální vzdálenost mezi úrovněmi
\def\SubtreeLabelOffset{0.30} % svislá mezera mezi trojúhelníkem a popiskem
\def\TriW{30pt}         % šířka trojúhelníku (podstromu)
\def\TriH{55pt}         % výška trojúhelníku (podstromu)

% ============================================================
% LEVÝ STROM  (kořen y, levý potomek x, pravý podstrom C)
% ============================================================

\node (Ly) at (0,0) {$y$};

\node (Lx) at (-\SiblingDist,-\LevelDist) {$x$};

\node[
    isosceles triangle,
    shape border rotate=90,
    anchor=apex,
    draw,
    fill=white,
    minimum height=\TriH,
    minimum width=\TriW,
    inner sep=0pt
] (LC) at ({0.9*\SiblingDist},-{1.15*\LevelDist}) {$C$};

\node[
    isosceles triangle,
    shape border rotate=90,
    anchor=apex,
    draw,
    fill=white,
    minimum height=\TriH,
    minimum width=\TriW,
    inner sep=0pt
] (LA) at ({-\SiblingDist-0.75*\SiblingDist},-{2.15*\LevelDist}) {$A$};

\node[
    isosceles triangle,
    shape border rotate=90,
    anchor=apex,
    draw,
    fill=white,
    minimum height=\TriH,
    minimum width=\TriW,
    inner sep=0pt
] (LB) at ({-\SiblingDist+0.75*\SiblingDist},-{2.15*\LevelDist}) {$B$};

\draw (Ly) -- (Lx);
\draw (Ly) -- (LC.apex);
\draw (Lx) -- (LA.apex);
\draw (Lx) -- (LB.apex);




% ============================================================
% PRAVÝ STROM  (kořen x, levý podstrom A, pravý potomek y)
% ============================================================

\def\XShift{6.5}

\node (Rx) at (\XShift,0) {$x$};

\node[
    isosceles triangle,
    shape border rotate=90,
    anchor=apex,
    draw,
    fill=white,
    minimum height=\TriH,
    minimum width=\TriW,
    inner sep=0pt
] (RA) at ({\XShift-0.9*\SiblingDist},-{1.15*\LevelDist}) {$A$};

\node (Ry) at ({\XShift+\SiblingDist},-\LevelDist) {$y$};

\node[
    isosceles triangle,
    shape border rotate=90,
    anchor=apex,
    draw,
    fill=white,
    minimum height=\TriH,
    minimum width=\TriW,
    inner sep=0pt
] (RB) at ({\XShift+\SiblingDist-0.75*\SiblingDist},-{2.15*\LevelDist}) {$B$};

\node[
    isosceles triangle,
    shape border rotate=90,
    anchor=apex,
    draw,
    fill=white,
    minimum height=\TriH,
    minimum width=\TriW,
    inner sep=0pt
] (RC) at ({\XShift+\SiblingDist+0.75*\SiblingDist},-{2.15*\LevelDist}) {$C$};

\draw (Rx) -- (RA.apex);
\draw (Rx) -- (Ry);
\draw (Ry) -- (RB.apex);
\draw (Ry) -- (RC.apex);




% ============================================================
% ŠIPKY ROTACÍ + POPISKY
% ============================================================

% velké šipky mezi kořeny (Ly = y vlevo) a (Rx = x vpravo)
% pomocné body pro stabilní centrování popisků
\path (Ly.east) ++(0,2mm) coordinate (RTopL);
\path (Rx.west) ++(0,2mm) coordinate (RTopR);
\path (Ly.east) ++(0,-2mm) coordinate (RBotL);
\path (Rx.west) ++(0,-2mm) coordinate (RBotR);

\draw[thick, -{Stealth[length=5mm,width=3.2mm]}]
    (RTopL) to[bend left=18] (RTopR);
\path (RTopL) -- (RTopR)
    node[midway, above=3mm,
         shape=rectangle, draw=none, fill=white,
         inner sep=1pt, font=\normalsize,
         text height=1.7ex, text depth=.3ex] {Pravá rotace};

\draw[thick, -{Stealth[length=5mm,width=3.2mm]}]
    (RBotR) to[bend left=18] (RBotL);
\path (RBotL) -- (RBotR)
    node[midway, below=3mm,
         shape=rectangle, draw=none, fill=white,
         inner sep=1pt, font=\normalsize,
         text height=1.7ex, text depth=.3ex] {Levá rotace};

\end{tikzpicture}
\caption{AVL rotace: pravá rotace (y,x,A,B,C) a levá rotace.}
\end{figure}